\section*{Bab \ref{ch:g10:synthesis-yield} --- Sintesis: Rendemen dan Kemurnian}

\begin{solution}{exo:g10:synthesis-yield:1}
$\eta = 4.2 / 5.0 = 0.84$, \omterm{def:g10:synthesis-yield:yield}{rendemen} \qty{84}{\percent}.
\end{solution}

\begin{solution}{exo:g10:synthesis-yield:2}
Menimbang \omterm{def:g7:chemical-reactions:reactant}{pereaksi}; \omterm{def:g10:synthesis-yield:reflux}{pemanasan dengan refluks}; \omterm{def:g10:synthesis-yield:vacuum-filtration}{penyaringan vakum} \omterm{def:g10:synthesis-yield:synthesis}{produk
kasar}; \omterm{def:g10:synthesis-yield:recrystallisation}{rekristalisasi}; mengukur titik leleh.
\end{solution}

\begin{solution}{exo:g10:synthesis-yield:3}
Jika dialirkan dari bawah, air memenuhi seluruh selubung sebelum keluar
di atas, sehingga seluruh panjang kondensor didinginkan. Ujung atasnya
dibiarkan terbuka agar peralatan itu tidak tertutup: memanaskan
peralatan yang tertutup akan menimbulkan tekanan dan dapat membuatnya
meledak.
\end{solution}

\begin{solution}{exo:g10:synthesis-yield:4}
Produk itu harus sangat mudah \omterm{def:g2:mixing-and-dissolving:dissolve}{larut} dalam \omterm{def:g6:solutions-and-solubility:solute}{pelarut} panas dan hampir tidak
\omterm{def:g2:mixing-and-dissolving:dissolve}{larut} dalam \omterm{def:g6:solutions-and-solubility:solute}{pelarut} dingin. Pengotornya, yang hanya sedikit, tetap
terlarut dalam \omterm{def:g6:solutions-and-solubility:solute}{pelarut} dingin dan berakhir di dalam \omterm{def:g3:separating-mixtures:filter}{filtrat}.
\end{solution}

\begin{solution}{exo:g10:synthesis-yield:5}
Tidak: sampel itu meleleh di bawah \qty{135}{\celsius} dan dalam rentang
beberapa derajat, tanda padatan yang tidak murni. Sampel itu harus
direkristalisasi, dikeringkan, lalu titik lelehnya diukur lagi.
\end{solution}

\begin{solution}{exo:g10:synthesis-yield:6}
$n = 2.00 / 138.0 = \qty{0.0145}{mol}$ asam salisilat, \omterm{def:g10:reaction-progress-table:limiting}{pereaksi
pembatas}; $m_{\max} = 0.0145 \times 180.0 = \qty{2.61}{g}$;
$\eta = 2.03 / 2.61 = 0.78$, yaitu \qty{78}{\percent}.
\end{solution}

\begin{solution}{exo:g10:synthesis-yield:7}
Massa yang ditimbang mencakup air, sehingga lebih besar daripada massa
aspirinnya: \omterm{def:g10:synthesis-yield:yield}{rendemen} yang dihitung terlalu tinggi, di sini bahkan
mustahil di atas \qty{100}{\percent}. \omterm{def:g10:synthesis-yield:synthesis}{Produk kasar} yang kering masih
mengandung pengotor (asam salisilat yang tidak bereaksi, produk
samping), yang juga menambah massa: \omterm{def:g10:synthesis-yield:yield}{rendemennya} juga terlalu tinggi
dibandingkan \omterm{def:g10:synthesis-yield:yield}{rendemen} sebenarnya aspirin murni.
\end{solution}

\begin{solution}{exo:g10:synthesis-yield:8}
\omterm{def:g10:synthesis-yield:vacuum-filtration}{Penyaringan vakum} jauh lebih cepat, dan meninggalkan kristal yang jauh
lebih kering, karena pengisapan menarik keluar sebagian besar cairan;
kristal-kristal itu juga mudah dicuci di atas corong dan dikerok dari
kertas saring yang datar.
\end{solution}

\begin{solution}{exo:g10:synthesis-yield:9}
Produk $3.1/5.0 = 0.62$; aspirin $3.1/5.0 = 0.62$; asam salisilat
$2.2/5.0 = 0.44$. Produk itu bergerak seperti aspirin dan tidak
memperlihatkan bercak asam salisilat: produk itu aspirin, tanpa asam
salisilat yang terdeteksi.
\end{solution}

\begin{solution}{exo:g10:synthesis-yield:10}
\omterm{def:g6:solutions-and-solubility:solute}{Pelarut} A: sangat mudah melarutkan dalam keadaan panas, hampir tidak
melarutkan dalam keadaan dingin; \omterm{def:g6:solutions-and-solubility:solute}{pelarut} B tetap melarutkan sebagian
besar produk bahkan dalam keadaan dingin. Dengan A, \qty{3.0}{g}
memerlukan sekitar $3.0 / 150 = \qty{0.020}{L}$, yaitu \qty{20}{mL}
\omterm{def:g6:solutions-and-solubility:solute}{pelarut} panas; setelah dingin, paling banyak
$2 \times 0.020 = \qty{0.04}{g}$ tetap terlarut.
\end{solution}

\begin{solution}{exo:g10:synthesis-yield:11}
Cairan yang mendidih memerlukan ruang di atasnya: buih dan percikan
tidak boleh mencapai kondensor. Batu didih memberi uap tempat-tempat
kecil untuk membentuk gelembung, sehingga cairannya mendidih dengan
tenang alih-alih menyembur tiba-tiba.
\end{solution}

\begin{solution}{exo:g10:synthesis-yield:12}
$0.80 \times 0.70 = 0.56$: \qty{56}{\percent}. Tiga langkah
\qty{90}{\percent}: $0.90^3 = 0.73$, \qty{73}{\percent}. Setiap langkah
kehilangan produk (juga waktu dan uang), dan kehilangan itu saling
berlipat: makin sedikit langkahnya, makin banyak produk yang tersisa di
akhir.
\end{solution}

\begin{solution}{exo:g10:synthesis-yield:13}
Paling banyak $3.3 \times 0.050 = \qty{0.17}{g}$ aspirin. Air yang
sangat dingin melarutkan aspirin lebih sedikit lagi, karena
\omterm{def:g6:solutions-and-solubility:solubility}{kelarutannya} turun bersama suhu.
\end{solution}

\begin{solution}{exo:g10:synthesis-yield:14}
\begin{enumerate}
  \item $n_0(\text{asam salisilat}) = 2.76 / 138.0 = \qty{0.0200}{mol}$,
        anhidrida \qty{0.0500}{mol}: asam salisilat adalah pembatas,
        $x_{\max} = \qty{0.0200}{mol}$, dan $0.0500 - 0.0200 =
        \qty{0.0300}{mol}$ anhidrida tersisa.
  \item \omterm{def:g10:synthesis-yield:synthesis}{Sintesis} itu menghasilkan \qty{0.0200}{mol} asam etanoat,
        penghancuran anhidrida yang berlebih
        $2 \times 0.0300 = \qty{0.0600}{mol}$: seluruhnya
        \qty{0.0800}{mol}.
\end{enumerate}
\end{solution}

\begin{solution}{exo:g10:synthesis-yield:15}
$n(\text{aspirin}) = \num{1.00e6} / 180.0 = \qty{5.56e3}{mol}$. Dengan
\omterm{def:g10:synthesis-yield:yield}{rendemen} \qty{85}{\percent}, asam salisilat yang diperlukan:
$\num{5.56e3} / 0.85 = \qty{6.54e3}{mol}$, yaitu
$\num{6.54e3} \times 138.0 = \qty{9.0e5}{g}$, sekitar \qty{0.90}{t}.
\end{solution}

\begin{solution}{pb:g10:synthesis-yield:1}
\textbf{1.} C: $7 + 4 = 11$ dan $9 + 2 = 11$; H: $6 + 6 = 12$ dan
$8 + 4 = 12$; O: $3 + 3 = 6$ dan $4 + 2 = 6$. Setara.

\textbf{2.} Asam salisilat $7 \times 12.0 + 6 \times 1.0 + 3 \times 16.0
= \qty{138.0}{g/mol}$; anhidrida $4 \times 12.0 + 6 \times 1.0 +
3 \times 16.0 = \qty{102.0}{g/mol}$; aspirin $9 \times 12.0 + 8 \times
1.0 + 4 \times 16.0 = \qty{180.0}{g/mol}$; asam etanoat $2 \times 12.0
+ 4 \times 1.0 + 2 \times 16.0 = \qty{60.0}{g/mol}$.

\textbf{3.} \omterm{def:g7:chemical-reactions:reactant}{Pereaksi}: asam salisilat dan anhidrida etanoat; produk yang
diinginkan: aspirin; produk samping: asam etanoat.

\textbf{4.} Pemanasan mempercepat reaksi; dengan refluks, uapnya
diembunkan dan dikembalikan, sehingga tidak ada \omterm{def:g7:chemical-reactions:reactant}{pereaksi} atau produk
yang hilang.

\textbf{5.} $n = 3.0 / 138.0 = \qty{0.0217}{mol}$.

\textbf{6.} $m = 6.0 \times 1.08 = \qty{6.48}{g}$;
$n = 6.48 / 102.0 = \qty{0.0635}{mol}$.

\textbf{7.} Asam salisilat: $0.0217 - x$; anhidrida $0.0635 - x$;
aspirin dan asam etanoat: $x$. Asam salisilat habis lebih dulu:
$x_{\max} = \qty{0.0217}{mol}$, asam salisilat adalah pembatas.

\textbf{8.} $m_{\max} = 0.0217 \times 180.0 = \qty{3.91}{g}$.

\textbf{9.} Agar semua asam salisilat diubah menjadi aspirin: sisanya
akan tetap tercampur dengan aspirin dan sukar dihilangkan, sedangkan
anhidrida yang berlebih dihancurkan oleh air yang ditambahkan pada
akhirnya, dan asam etanoat yang dihasilkannya tercuci.

\textbf{10.} Massa itu mencakup air: \qty{3.6}{g} akan memberikan
\omterm{def:g10:synthesis-yield:yield}{rendemen} palsu $3.6 / 3.91 = \qty{92}{\percent}$.

\textbf{11.} $3.3 / 3.91 = 0.84$: \qty{84}{\percent} \omterm{def:g10:synthesis-yield:synthesis}{produk kasar}
kering, yang tidak murni.

\textbf{12.} Pengotornya dihilangkan, dan sebagian aspirin tetap
terlarut dalam \omterm{def:g6:solutions-and-solubility:solute}{pelarut} dingin dan dalam air cucian.

\textbf{13.} Paling banyak $3.3 \times 0.030 = \qty{0.10}{g}$.

\textbf{14.} $\eta = 2.7 / 3.91 = 0.69$: \qty{69}{\percent}.

\textbf{15.} Titik leleh yang tajam pada nilai aspirin murni: kristal
itu aspirin, dan murni.

\textbf{16.} Meleleh di bawah \qty{135}{\celsius} dan dalam rentang yang
lebar: \omterm{def:g10:synthesis-yield:synthesis}{produk kasarnya} tidak murni.

\textbf{17.} \omterm{def:g10:synthesis-yield:synthesis}{Produk kasar} mengandung aspirin dan asam salisilat yang
tidak bereaksi; \omterm{def:g10:synthesis-yield:recrystallisation}{rekristalisasi} telah menghilangkan asam salisilat itu.

\textbf{18.} $R_f = 3.1 / 5.0 = 0.62$.

\textbf{19.} Kira-kira sama banyaknya: dari \qty{3.91}{g} yang mungkin
menjadi \qty{3.3}{g} \omterm{def:g10:synthesis-yield:synthesis}{produk kasar}, dan dari \qty{3.3}{g} menjadi
\qty{2.7}{g} pada \omterm{def:g10:synthesis-yield:recrystallisation}{rekristalisasi}, masing-masing sekitar \qty{0.6}{g}.
Karena \omterm{def:g10:synthesis-yield:synthesis}{produk kasarnya} bukan aspirin murni, bagian pertama pekerjaan itu
sebenarnya kehilangan aspirin sedikit lebih banyak daripada
\qty{0.6}{g}.

\textbf{20.} \omterm{def:g10:synthesis-yield:yield}{Rendemen} \omterm{def:g10:synthesis-yield:synthesis}{sintesis} ini sekitar \qty{69}{\percent}.
\end{solution}
